\chapter*{Abreviaturas y símbolos}
\addcontentsline{toc}{chapter}{Abreviaturas y símbolos}

\begin{longtable}{p{0.20\textwidth}p{0.72\textwidth}}
\toprule
\textbf{Abreviatura} & \textbf{Significado} \\
\midrule
BOLD & Blood oxygenation level dependent; contraste dependiente de oxigenación sanguínea. \\
fMRI & Resonancia magnética funcional. \\
GLM & Modelo lineal general. \\
HCP & Human Connectome Project. \\
HRF & Función de respuesta hemodinámica. \\
M1 & Corteza motora primaria. \\
MLP & Perceptrón multicapa. \\
ODE & Ecuación diferencial ordinaria. \\
PINN & Red neuronal informada por la física. \\
ROI & Región de interés. \\
SNR & Relación señal--ruido. \\
\bottomrule
\end{longtable}

\begin{longtable}{p{0.20\textwidth}p{0.72\textwidth}}
\toprule
\textbf{Símbolo} & \textbf{Interpretación en el modelo} \\
\midrule
$u(t)$ & Entrada neuronal o estímulo experimental. \\
$s(t)$ & Señal vasodilatadora. \\
$f(t)$ & Flujo sanguíneo normalizado respecto del basal. \\
$v(t)$ & Volumen sanguíneo venoso normalizado. \\
$q(t)$ & Contenido normalizado de desoxihemoglobina. \\
$y(t)$ & Cambio BOLD fraccional. \\
$\epsilon$ & Eficacia de la entrada neuronal sobre la señal vasodilatadora. \\
$\kappa_s$ & Tasa de decaimiento de la señal vasodilatadora. \\
$\kappa_f$ & Retroalimentación autorreguladora del flujo. \\
$\tau$ & Tiempo medio de tránsito sanguíneo. \\
$\alpha$ & Exponente de la relación flujo--volumen. \\
$E_0$ & Fracción de extracción de oxígeno en reposo. \\
$V_0$ & Fracción de volumen sanguíneo venoso basal. \\
$\lambda_d,\lambda_p,\lambda_i,\lambda_t$ & Pesos de datos, física, condición inicial y condición terminal. \\
\bottomrule
\end{longtable}
